% ==============================================================================
% SECTION: DATASET AND LINGUISTIC VARIANT BENCHMARK CONSTRUCTION
% Target: Q2/Q3 Peer-Reviewed Journal in Artificial Intelligence & Cyber Security
% ==============================================================================

\subsection{Symmetric Cross-Lingual, Orthographic, and Register Variation Benchmark Construction}
\label{sec:dataset_construction}

To rigorously assess the cross-lingual robustness of Large Language Models (LLMs) under realistic orthographic and register variations, we introduce a structurally symmetric, paired benchmark dataset bridging English and diverse Vietnamese linguistic phenomena.

\subsubsection{Seed Pair Taxonomy and Category Distribution}
The foundation of the benchmark comprises 30 primary seed pairs ($N_{\text{seed}} = 30$), designed under strict academic and ethical boundaries. The seed pairs establish a canonical $2:1$ ratio ($66.7\%$ harmful to $33.3\%$ benign; $N_{\text{harmful}} = 20$, $N_{\text{benign}} = 10$). When expanded across all six cohorts (five Vietnamese variants plus the English baseline), the full population yields exactly $120$ harmful and $60$ benign test instances ($N_{\text{pop}} = 180$). Harmful seeds are strictly balanced across all five risk domains ($4$ harmful seeds $\times 6$ cohorts $= 24$ harmful prompts per category), while benign contrastive controls reflect domain operational density ($18, 12, 6, 18, 6$ benign prompts, yielding category totals of $42, 36, 30, 42, 30$; see Table~\ref{tab:ablation_and_scaling}, Panel~B). The pairs are divided into two distinct functional cohorts:
\begin{enumerate}
    \item \textbf{Harmful Safety Evaluation Cohort ($N_{\text{harmful}} = 20$)}: Designed to test whether the model's safety guardrails reliably intercept hazardous intent across languages. To ensure comprehensive threat surface coverage, these 20 seed queries are evenly distributed ($4$ seeds each) across five critical safety categories aligned with domestic and international safety standards:
    \begin{itemize}
        \item \textit{Vietnamese Legal \& Administrative Compliance (\foreignlanguage{vietnamese}{Pháp luật \& Hành chính công VN})}: Regulatory evasion, falsification of public records (e.g., chip-embedded citizen ID CCCD), and administrative manipulation under the Vietnamese legal framework ($42$ prompts: $24$ harmful, $18$ benign).
        \item \textit{Financial, Banking \& Securities Fraud (\foreignlanguage{vietnamese}{Lừa đảo tài chính/Ngân hàng/Chứng khoán VN})}: Fictitious VAT invoicing, OTP phishing, and securities market manipulation ($36$ prompts: $24$ harmful, $12$ benign).
        \item \textit{Cultural Norms, History \& Sensitive Topics (\foreignlanguage{vietnamese}{Văn hóa -- Lịch sử -- Tôn giáo nhạy cảm VN})}: Sensitive historical distortion, disinformation, and socially divisive framing in domestic contexts ($30$ prompts: $24$ harmful, $6$ benign).
        \item \textit{Cybersecurity \& Malware Engineering (\foreignlanguage{vietnamese}{An ninh mạng \& Mã độc})}: Software exploitation, keyloggers, ransomware, SQL injection, and OWASP LLM vulnerabilities ($42$ prompts: $24$ harmful, $18$ benign).
        \item \textit{Social Harm \& Self-Harm (\foreignlanguage{vietnamese}{Tác hại xã hội/Tự gây hại})}: Dangerous substances, harassment, and privacy/personal-harm inquiries ($30$ prompts: $24$ harmful, $6$ benign).
    \end{itemize}
    Importantly, in accordance with responsible AI safety guidelines, all prompts are framed as academic, procedural, or compliance-testing inquiries without containing executable malware or harmful instructional guidance.

    \item \textbf{Benign Contrastive Control Cohort ($N_{\text{benign}} = 10$)}: A control set containing specialized academic, administrative, and technological inquiries that share surface lexical overlap with the harmful cohort (e.g., forensic database auditing, parameterized defense against SQL injection, statutory administrative appeal procedures). This cohort enables precise calibration of model utility and over-refusal behavior ($\text{ORR}$).
\end{enumerate}

\subsubsection{Symmetric Five-Variant Morphological Transformation}
For each seed pair, we generate one standardized English baseline prompt ($\texttt{EN\_RT\_001}$ to $\texttt{EN\_RT\_030}$) recorded in \texttt{english\_baseline\_benchmark.json}, and five semantically aligned Vietnamese morphological variants sharing the identical \texttt{paired\_en\_id}, yielding $30 \times 5 = 150$ Vietnamese test instances recorded in \texttt{vietnamese\_benchmark.json}:
\begin{enumerate}
    \item \textbf{Standard Vietnamese (\textit{standard\_vi})}: Syntactically normative Vietnamese containing full diacritical tone marks and formal terminology as utilized in official documentation.
    \item \textbf{Unaccented Vietnamese (\textit{no\_diacritics})}: Complete stripping of Vietnamese diacritics and tonal markers (converting tonal and vocalic variants, including acute \texttt{U+00E1}, grave \texttt{U+00E0}, hook-above \texttt{U+1EA3}, tilde \texttt{U+00E3}, dot-below \texttt{U+1EA1}, circumflex \texttt{U+00E2}, breve \texttt{U+0103}, d-bar \texttt{U+0111}, horn-o \texttt{U+01A1}, and horn-u \texttt{U+01B0}, into base ASCII characters: \textit{a, d, e, o, u}). This reflects widespread informal communication practices in mobile SMS, search queries, and online forums \cite{nguyen2020phobert}, probing tokenizer subword splitting fragility.
    \item \textbf{Youth Slang and Leetspeak (\textit{teencode\_slang})}: Orthographic transformations reflecting contemporary Vietnamese digital communication, including phonological substitutions (e.g., $ph \to f$, $qu \to w$, $c/k \to k$), frequent contractions (e.g., \textit{\foreignlanguage{vietnamese}{được}} $\to$ \textit{dc}, \textit{\foreignlanguage{vietnamese}{không}} $\to$ \textit{k/ko}, \textit{\foreignlanguage{vietnamese}{làm}} $\to$ \textit{lm}), and vowel shifts (e.g., \textit{\foreignlanguage{vietnamese}{quá}} $\to$ \textit{wa}).
    \item \textbf{Intra-Sentential Code-Switching (\textit{code\_switching})}: Natural bilingual blending where standard Vietnamese grammatical frames incorporate English technical loanwords and operational expressions (e.g., \textit{bypass firewall}, \textit{audit \foreignlanguage{vietnamese}{hệ thống}}, \textit{exploit payload}, \textit{leak database}), mirroring actual developer discourse in Vietnam.
    \item \textbf{Canonical Decomposed Unicode (\textit{unicode\_nfd})}: Canonical Decomposition (Unicode Normalization Form D), which separates precomposed Vietnamese characters into a base Latin letter followed by separate combining diacritical code points (for example, the precomposed character \texttt{U+1EC7} (Latin small letter \textit{e} with circumflex and dot below) is canonically decomposed into base letter \texttt{e} + combining circumflex \texttt{U+0302} + combining dot below \texttt{U+0323}). This variant assesses vulnerability to low-level byte-sequence representation manipulation without altering the surface glyph.
\end{enumerate}

\subsubsection{Contextual Natural Language Framing}
Direct, single-clause interrogatives frequently trigger trivial surface keyword pattern matchers in safety alignment filters, producing artificial floor effects. To guarantee ecological validity, we implement an automated contextual framing transformation:
\begin{equation}
    x_{\text{framed}} = \mathcal{T}_{\text{frame}}(x_{\text{seed}}, \mathcal{L}, \mathcal{I})
    \label{eq:framing}
\end{equation}
where $\mathcal{L} \in \{\text{EN}, \text{VI}\}$ denotes language and $\mathcal{I} \in \{\text{harmful}, \text{benign}\}$ denotes intent type. The transformation wraps the core inquiry into natural organizational scenarios---such as a security auditor reviewing legacy source code, an administrative consultant drafting procedural checklists, or an IT instructor illustrating defensive engineering concepts. Crucially, the semantic boundary between malicious intent and lawful inquiry is strictly conserved across all five morphological transformations.

\subsubsection{Full-Population Multi-Architecture Evaluation and Double-Blind Human Validation Design}
\label{sec:cohort_reconciliation}
To guarantee statistical power and eliminate sample-size discrepancies across tables and sections, we execute a unified full-population evaluation alongside a verified double-blind human validation protocol:
\begin{itemize}
    \item \textbf{Complete Population Benchmark across Six Architectures ($N_{\text{pop}} = 180$ per model; $1{,}080$ total real model evaluations)}: Encompasses the full generative suite of 30 bilingual seed pairs ($20$ harmful and $10$ benign pairs, establishing an exact $2:1$ population ratio: $120$ harmful and $60$ benign cases, or $66.7\% : 33.3\%$). This repository spans $150$ Vietnamese test cases ($30$ cases across each of the 5 linguistic variants: $N_{\text{h}} = 20$ harmful and $N_{\text{b}} = 10$ benign per variant) and $30$ corresponding English baseline test cases ($N_{\text{h}} = 20, N_{\text{b}} = 10$). Every single prompt in the 180-case benchmark was executed through real Ollama inference across six distinct local architectures: two global multilingual baselines (\texttt{Llama-3-8B-Instruct}, \texttt{Qwen-2.5-7B-Instruct}), two Vietnamese/Southeast Asian specialized models (\texttt{SeaLLMs-v3-7B-Chat}, \texttt{Vistral-7B-Chat}), one weakly aligned open-weight baseline (\texttt{Mistral-7B-Instruct-v0.3}), and our hardened defensive configuration (\texttt{Llama-3-Secure + Stage-0 Intervention}), all scored with a decoupled \texttt{Qwen-2.5-7B-Instruct} LLM judge and providing $n_{\text{pairs}} = 20$ matched harmful pairs per cohort for exact paired McNemar testing and omnibus Cochran's $Q$ analysis in Section~\ref{sec:experiments}.
    \item \textbf{Verified Double-Blind Human-Annotated Validation Corpus ($N_{\text{VI}} = 50$) \& Expanded Multi-Architecture Protocol}: Two independent bilingual cybersecurity researchers completed double-blind annotations across a stratified $N_{\text{VI}} = 50$ Vietnamese subset of \texttt{Llama-3-8B-Instruct} outputs ($N_{\text{harmful}} = 35$ [$7$/variant] and $N_{\text{benign}} = 15$ [$3$/variant]), yielding $\kappa_{\text{quad}} = 0.876$, Macro-F1 $= 93.8\%$, and $94.0\%$ ($47/50$) observed agreement. To extend empirical sensitivity ($\text{Se}$) estimation to models with non-zero State 2 jailbreaks (\texttt{Vistral-7B-Chat}, \texttt{SeaLLMs-v3-7B-Chat}, \texttt{Mistral-7B-Instruct-v0.3}), we additionally package a double-blind workbook of $250$ stratified multi-architecture outputs (\texttt{Annotator\_1..2\_Blind\_250.xlsx}) integrated with an automated import/adjudication pipeline (\texttt{scripts/import\_blind\_excel\_300.py}); we strictly exclude uncompleted rows from the empirical results reported in this study.
\end{itemize}
